\section{Méthodes}
\noindent

\subsection{Récupération et préparations des données}
\noindent

\subsubsection{Sources et périmètre}


\begin{table}[H]
    \centering
    \fontsize{10}{12}\selectfont
    \renewcommand{\arraystretch}{1.3}
    \caption{Source des données utilisées pour l'analyse et la modélisation.}
    \begin{tabularx}{\linewidth}{p{2.5cm} X X X X X}
        \toprule
        Données & Source & Période \newline disponible & Granularité initiale & Variables \newline retenues & Unités \\
        \midrule
        Consommation \newline électrique & \cite{rte_eco2mix_2026} & Depuis 2012 & 30 minutes & Consommation & MWh \\
        Météo & \cite{meteo_france_synop_2026} & Depuis 1996 & 3 heures & Température, précipitations & ° K, mm \\
        Calendrier & N/A & N/A & 1 jour & Date, jours \newline fériés et \newline vacances & N/A \\
        \bottomrule
    \end{tabularx}
    
    \label{tab:scenarios}
\end{table}

\subsubsection{Construction du tableau horaire}

- Consommation : nettoyage, agrégation des demi-heures et convention retenue pour dater la moyenne. \newline
- Météo : stations métropolitaines retenues, seuil de couverture de la température et des précipitations, moyenne spatiale à poids égaux, conversion des unités et reconstruction horaire. \newline
- Préciser que l'écart-type météo mesure la dispersion entre stations, pas l'incertitude de la moyenne nationale.\newline
- Calendrier : jour, heure, week-end, jours fériés et vacances par zone ; justification des variables retenues.\newline
- Jointure temporelle, vérification du fuseau RTE, traitement des doublons, valeurs manquantes et observations atypiques.\newline
- Changements d'heure : expliquer la convention permettant de traiter la demande de 24 cibles malgré les journées françaises de 23 ou 25 heures ; documenter toute exclusion.

\subsubsection{Analyse exploratoire}

Explorer les données d'apprentissage pour orienter les modèles sans utiliser le test final pour choisir les variables.

- Cycles journalier et hebdomadaire, saisonnalité annuelle. \newline
- Relation température–consommation et effets calendaires.\newline
- Tendances, ruptures de régime et observations atypiques, sans suppression automatique des événements réels.

Figures proposées : profils horaires et saisonniers ; relation température–consommation ou de série longue montrant les ruptures. 


\subsection{Benchmarks et modèles}
\noindent

\begin{algorithm}[h]
    \caption{Un superbe algorithme}\label{alg:justanalgorithm}
    \textit{Entrées}
    \begin{itemize}
        
        \item[$\bullet$] $Variable = $ description de la variable
        \item[$\bullet$] $Variable = $ description de la variable
        \item[$\bullet$] $Variable = $ description de la variable
        \item[$\bullet$] $Variable = $ description de la variable
    \end{itemize}
    \textit{Initialisations}
    \begin{itemize}
        \item[$\bullet$] $Variable = valeur$
        \item[$\bullet$] Effectuer une action
    \end{itemize}
    \textit{Calculs}\\
        \For{$variable \gets 0$ \KwTo $Variable$ \KwBy $1$}{
            \ForEach{$truc \in machin$}{
                Faire à $truc$ quelque chose prenant en compte $Variable$ \\
                Mise à jour de $machin$
            }
        $Variable = Variable \times Variable$
        }
    \textit{Sorties}
    \begin{itemize}
        \item[$\bullet$] Ce que l'on sort de l'algorithme
    \end{itemize}
\end{algorithm}

\subsection{Validation chronologique et métriques}
\noindent

\subsubsection{Sous-sous-section}
\noindent
